
\documentclass{article}
\usepackage{colortbl}
\usepackage{makecell}
\usepackage{multirow}
\usepackage{supertabular}

\begin{document}

\newcounter{utterance}

\centering \large Interaction Transcript for game `codenames', experiment `risk\_low', episode 1 with reflexion{-}Qwen3.5{-}27B.
\vspace{24pt}

{ \footnotesize  \setcounter{utterance}{1}
\setlength{\tabcolsep}{0pt}
\begin{supertabular}{c@{$\;$}|p{.15\linewidth}@{}p{.15\linewidth}p{.15\linewidth}p{.15\linewidth}p{.15\linewidth}p{.15\linewidth}}
    \# & $\;$A & \multicolumn{4}{c}{Game Master} & $\;\:$B\\
    \hline

    \theutterance \stepcounter{utterance}  
    & & & \multicolumn{2}{p{0.3\linewidth}}{
        \cellcolor[rgb]{0.9,0.9,0.9}{
            \makecell[{{p{\linewidth}}}]{
                \texttt{\tiny{[GM$|$GM]}}
                {'hidden': {'team': ['fish'], 'innocent': ['cell', 'cycle', 'night'], 'opponent': [], 'assassin': []}, 'revealed': {'team': {'team': [], 'innocent': [], 'opponent': [], 'assassin': []}, 'opponent': {'team': [], 'innocent': [], 'opponent': [], 'assassin': []}}}
            }
        }
    }
    & & \\ \\

    \theutterance \stepcounter{utterance}  
    & & \multicolumn{4}{p{0.6\linewidth}}{
        \cellcolor[rgb]{0.9,0.9,0.9}{
            \makecell[{{p{\linewidth}}}]{
                \texttt{\tiny{[P1$\langle$GM]}}
                \texttt{You have attempted this same task before and did not succeed. These are the plans you wrote after those attempts.} \\
\texttt{Trial \#1: New plan: My previous strategy failed because I consistently provided clues that were too broad, causing my partner to guess distractor words (like "wave," "organ," "cross," "theater," and "telescope") which ended my turns prematurely. I attempted to guess multiple words at once (e.g., "band" and "game" with "play"), but the ambiguity led to errors. Next time, I will prioritize safety over speed by selecting clues that are highly specific to a single target word or a very tight semantic cluster that excludes all distractors. For example, instead of "music" (which triggered "organ"), I should have used "drums" or "guitar" for "band." Instead of "snake" (which triggered "cross"), I should have used "reptile" or "venom" for "fish" if grouping, or simply "aquatic" for "fish" alone. I will avoid clues that have strong associations with the remaining distractor list and, if a clue risks a distractor, I will limit the TARGETS list to only the most obvious match to prevent the partner from guessing incorrectly.} \\
\\ 
\texttt{Trial \#2: New plan: My previous attempt failed because I provided a clue ("aquatic") that, while semantically correct for "fish," was too broad and likely triggered a distractor word (such as "cell" in a biological context or "cycle" in an ecological context) or simply failed to secure the win due to the high risk associated with the specific distractor list. In this specific instance with only one target word ("fish") and distractors "cell," "cycle," and "night," the clue "aquatic" was insufficiently distinct. Next time, I will select a clue that creates a stronger, more exclusive association with the target while explicitly avoiding the semantic fields of the distractors. For "fish," I should have used a more specific term like "scale," "fin," or "hook" which are directly related to the animal but have no logical connection to "cell," "cycle," or "night." I will ensure the clue is narrow enough to prevent the partner from guessing any of the listed distractors, prioritizing a safe, unambiguous link over a broad category.} \\
\\ 
\texttt{Trial \#3: New plan: My previous attempt failed because I selected "fin" as a clue for "fish," which, while semantically correct, was likely too specific or ambiguous in the context of the distractors ("cell," "cycle," "night"), causing the partner to miss the target or guess incorrectly. Although "fin" is a part of a fish, it might not be the most immediate association compared to other potential clues, or it may have inadvertently triggered a distractor if the partner overthought the connection. Next time, I will choose a clue that represents the most direct, high{-}frequency association with the target word while ensuring zero overlap with the distractors. For "fish," I should have used "swim" or "water," which are universally associated with fish but have no logical connection to "cell," "cycle," or "night." I will prioritize the most obvious, common{-}sense link to ensure the partner guesses the target immediately without hesitation or confusion.} \\
\\ 
\\ 
\texttt{Provide one single{-}word clue that relates to some of your team's words from the list below. You can choose to how many and to which words your clue relates to.} \\
\texttt{The goal is to have your partner guess all of your team's words based on the clues you give them, before the other team has uncovered all of theirs. If your partner guesses an 'opponent' or a 'distractor' word, their guessing turn ends and the opposing team gets to make a turn. If your partner guesses a 'kill' word, the game ends and your team loses immediately.} \\
\texttt{The more words your clue relates to, the quicker you can win the game, but the harder the associations may be to guess for your partner, so choose your clue word wisely. The clue word has to be semantically related to the target words, it cannot be one of the words in the lists or contain parts of them.} \\
\\ 
\texttt{Always give your single{-}word clue and your comma{-}separated list of related target words in the following format and make your answers as short as possible, never include any other text than is required in this form:} \\
\\ 
\texttt{CLUE: <WORD>} \\
\texttt{TARGETS: <TARGETS>} \\
\\ 
\texttt{Your team words are: fish.} \\
\texttt{Your opponent's team words are: .} \\
\texttt{Distractor words are: cell, cycle, night.} \\
\texttt{The kill words are: .} \\
            }
        }
    }
    & \\ \\

    \theutterance \stepcounter{utterance}  
    & \multicolumn{4}{p{0.6\linewidth}}{
        \cellcolor[rgb]{0.9,0.9,0.9}{
            \makecell[{{p{\linewidth}}}]{
                \texttt{\tiny{[P1$\rangle$GM]}}
                \texttt{CLUE: SWIM} \\
\texttt{TARGETS: fish} \\
            }
        }
    }
    & & \\ \\

    \theutterance \stepcounter{utterance}  
    & & & \multicolumn{2}{p{0.3\linewidth}}{
        \cellcolor[rgb]{0.9,0.9,0.9}{
            \makecell[{{p{\linewidth}}}]{
                \texttt{\tiny{[GM$|$GM]}}
                \texttt{swim} \\
            }
        }
    }
    & & \\ \\

    \theutterance \stepcounter{utterance}  
    & & & \multicolumn{2}{p{0.3\linewidth}}{
        \cellcolor[rgb]{0.9,0.9,0.9}{
            \makecell[{{p{\linewidth}}}]{
                \texttt{\tiny{[GM$|$GM]}}
                ['fish']
            }
        }
    }
    & & \\ \\

    \theutterance \stepcounter{utterance}  
    & & & \multicolumn{2}{p{0.3\linewidth}}{
        \cellcolor[rgb]{0.9,0.9,0.9}{
            \makecell[{{p{\linewidth}}}]{
                \texttt{\tiny{[GM$|$GM]}}
                {'word': 'fish', 'assignment': 'team'}
            }
        }
    }
    & & \\ \\

    \theutterance \stepcounter{utterance}  
    & & & \multicolumn{2}{p{0.3\linewidth}}{
        \cellcolor[rgb]{0.9,0.9,0.9}{
            \makecell[{{p{\linewidth}}}]{
                \texttt{\tiny{[GM$|$GM]}}
                \texttt{opponent has won} \\
            }
        }
    }
    & & \\ \\

    \theutterance \stepcounter{utterance}  
    & \multicolumn{4}{p{0.6\linewidth}}{
        \cellcolor[rgb]{0.95,0.95,0.95}{
            \makecell[{{p{\linewidth}}}]{
                \texttt{\tiny{[P2$\rightarrow$P1]}}
                \texttt{Post{-}episode Reflexion (not sent to GM):} \\
\texttt{New plan: My previous attempt failed because I chose "SWIM" as a clue for "fish," which, while semantically related, is an action that can be associated with the distractor "cycle" (e.g., swimming cycle, or the cycle of life in water) or simply failed to be the most distinct link given the specific distractor list. Although "SWIM" is a strong association, the game ended in a loss, suggesting the partner either guessed a distractor or the clue was deemed too risky. Next time, I will select a clue that is a concrete noun or a specific attribute of the target that has absolutely no semantic overlap with "cell," "cycle," or "night." For "fish," I should have used "scale," "fin," or "hook." These are physical parts or tools directly tied to fish but have no logical connection to the distractors. I will prioritize a clue that is a unique identifier for the target word to ensure the partner guesses "fish" immediately without considering the distractors.} \\
            }
        }
    }
    & & \\ \\

\end{supertabular}
}

\end{document}
